\documentclass[11pt]{article}

\usepackage[a4paper,margin=1in]{geometry}
\usepackage{amsmath,amssymb,amsthm,mathtools}
\usepackage[T1]{fontenc}
\usepackage{lmodern}
\usepackage{microtype}
\usepackage{hyperref}

\newtheorem{theorem}{Theorem}[section]
\newtheorem{lemma}[theorem]{Lemma}
\newtheorem{proposition}[theorem]{Proposition}
\newtheorem{corollary}[theorem]{Corollary}
\newtheorem{conjecture}[theorem]{Conjecture}
\newtheorem{remark}[theorem]{Remark}
\newtheorem{definition}[theorem]{Definition}

\title{The Connectedness Threshold Problem\\
for the Clean OBC Hatano--Nelson Family}
\author{}
\date{}

\begin{document}
\maketitle

\section{Setup}

We work with the real nonnormal tridiagonal Toeplitz matrices
\[
\mathbf{A}_n=\operatorname{tridiag}(a,0,b)
=\begin{pmatrix}
0 & b \\
a & 0 & \ddots \\
& \ddots & \ddots & \ddots \\
&& \ddots & 0 & b \\
&&& a & 0
\end{pmatrix}
\in\mathbb{R}^{n\times n},
\quad n>1, \quad 0<a<b.
\]

For $\varepsilon>0$, the $\varepsilon$-pseudospectrum of $\mathbf{A}_n$ is
\[
\sigma_\varepsilon(\mathbf{A}_n)\coloneqq\{z\in\mathbb{C}\colon\|(z\mathbf{I}-\mathbf{A}_n)^{-1}\|>\varepsilon^{-1}\},
\]
where $\|\cdot\|$ is the standard $2$-norm, with the convention that
$\|(z\mathbf{I}-\mathbf{A}_n)^{-1}\|=+\infty$ for $z\in\sigma(\mathbf{A}_n)$.
Equivalently,
\[
\sigma_\varepsilon(\mathbf{A}_n)
=\{z\in\mathbb{C}\colon s_{\min}(z\mathbf{I}-\mathbf{A}_n)<\varepsilon\},
\]
where $s_{\min}(\cdot)$ denotes the smallest singular value.

\section{An Algebraic Boundary Polynomial}

For $z=x+iy$, define
\[
\mathbf{H}_{n,\varepsilon}(x,y)
\coloneqq
\bigl((z\mathbf{I}-\mathbf{A}_n)^*(z\mathbf{I}-\mathbf{A}_n)-\varepsilon^2\mathbf{I}\bigr),
\]
and set
\[
F_{n,\varepsilon}(x,y)\coloneqq \det \mathbf{H}_{n,\varepsilon}(x,y).
\]
We shall call $F_{n,\varepsilon}$ the \emph{$\varepsilon$-pseudospectral boundary polynomial}
of $\mathbf{A}_n$.

\begin{proposition}\label{prop:boundary-polynomial}
For every $n>1$ and $\varepsilon>0$, the polynomial $F_{n,\varepsilon}(x,y)$ belongs to
$\mathbb{R}[x^2,y^2]$ and has total degree $2n$. Moreover,
\[
\partial\sigma_\varepsilon(\mathbf{A}_n)
=\bigl\{(x,y)\in\mathbb{R}^2\colon F_{n,\varepsilon}(x,y)=0,
\ \mathbf{H}_{n,\varepsilon}(x,y)\succeq 0\bigr\}
\subseteq
\bigl\{(x,y)\in\mathbb{R}^2\colon F_{n,\varepsilon}(x,y)=0\bigr\}.
\]
\end{proposition}

\begin{proof}
Since
\[
\sigma_\varepsilon(\mathbf{A}_n)
=\{z\in\mathbb{C}\colon s_{\min}(z\mathbf{I}-\mathbf{A}_n)<\varepsilon\},
\]
and the map $z\mapsto s_{\min}(z\mathbf{I}-\mathbf{A}_n)$ is continuous, its boundary is
\[
\partial\sigma_\varepsilon(\mathbf{A}_n)
=\{z\in\mathbb{C}\colon s_{\min}(z\mathbf{I}-\mathbf{A}_n)=\varepsilon\}.
\]
The eigenvalues of
\[
\mathbf{H}_{n,\varepsilon}(x,y)
=\bigl((z\mathbf{I}-\mathbf{A}_n)^*(z\mathbf{I}-\mathbf{A}_n)-\varepsilon^2\mathbf{I}\bigr)
\]
are precisely $s_j(z\mathbf{I}-\mathbf{A}_n)^2-\varepsilon^2$, where $s_j$ are the singular
values of $z\mathbf{I}-\mathbf{A}_n$. Hence
\[
s_{\min}(z\mathbf{I}-\mathbf{A}_n)=\varepsilon
\iff
\lambda_{\min}\bigl(\mathbf{H}_{n,\varepsilon}(x,y)\bigr)=0
\iff
\bigl(F_{n,\varepsilon}(x,y)=0\ \text{and}\ \mathbf{H}_{n,\varepsilon}(x,y)\succeq 0\bigr).
\]
This proves the boundary characterization.

It remains to prove that $F_{n,\varepsilon}\in\mathbb{R}[x^2,y^2]$ and that
$\deg F_{n,\varepsilon}=2n$. First, $\mathbf{A}_n$ is real, so replacing $y$ by $-y$
conjugates the entries of $\mathbf{H}_{n,\varepsilon}(x,y)$; therefore
\[
F_{n,\varepsilon}(x,-y)=F_{n,\varepsilon}(x,y).
\]
Next, if
\[
\mathbf{J}_n\coloneqq \operatorname{diag}(1,-1,1,-1,\dots,(-1)^{n-1}),
\]
then $\mathbf{J}_n\mathbf{A}_n\mathbf{J}_n^{-1}=-\mathbf{A}_n$. Hence
\[
\mathbf{J}_n\mathbf{H}_{n,\varepsilon}(x,y)\mathbf{J}_n^{-1}
=\bigl((z\mathbf{I}+\mathbf{A}_n)^*(z\mathbf{I}+\mathbf{A}_n)-\varepsilon^2\mathbf{I}\bigr),
\]
so
\[
F_{n,\varepsilon}(-x,-y)=F_{n,\varepsilon}(x,y).
\]
Combining this with the evenness in $y$ gives $F_{n,\varepsilon}(-x,y)=F_{n,\varepsilon}(x,y)$.
Thus $F_{n,\varepsilon}$ is even in both $x$ and $y$, that is,
$F_{n,\varepsilon}\in\mathbb{R}[x^2,y^2]$.

Finally, each diagonal entry of $\mathbf{H}_{n,\varepsilon}(x,y)$ has degree $2$ in $(x,y)$,
all first off-diagonal entries have degree $1$, and all second off-diagonal entries have degree $0$.
The product of the diagonal entries contributes the top-degree term
$(x^2+y^2)^n$, so $\deg F_{n,\varepsilon}=2n$.
\end{proof}

\begin{remark}\label{rem:not-exact-algebraic-curve}
In general, the full zero set $\{F_{n,\varepsilon}=0\}$ is larger than the true boundary
$\partial\sigma_\varepsilon(\mathbf{A}_n)$. The equation $F_{n,\varepsilon}(x,y)=0$ only says that
\emph{some} singular value of $z\mathbf{I}-\mathbf{A}_n$ equals $\varepsilon$; the boundary is the
branch on which the \emph{smallest} singular value equals $\varepsilon$.
\end{remark}

\begin{proposition}\label{prop:explicit-H}
Set
\[
s\coloneqq x^2+y^2-\varepsilon^2,
\qquad
c\coloneqq az+b\bar z=(a+b)x-i(b-a)y,
\qquad
p\coloneqq ab,
\qquad
d\coloneqq s+a^2+b^2.
\]
Then
\[
\mathbf{H}_{n,\varepsilon}(x,y)
=(|z|^2-\varepsilon^2)\mathbf{I}-\bar z\,\mathbf{A}_n-z\,\mathbf{A}_n^{\mathsf{T}}+\mathbf{A}_n^{\mathsf{T}}\mathbf{A}_n,
\]
and its entries are given by
\[
\bigl(\mathbf{H}_{n,\varepsilon}(x,y)\bigr)_{jk}
=
\begin{cases}
s+a^2, & j=k=1,\\
s+b^2, & j=k=n,\\
d, & j=k\in\{2,\dots,n-1\},\\
-c, & k=j+1,\\
-\bar c, & j=k+1,\\
p, & |j-k|=2,\\
0, & |j-k|>2.
\end{cases}
\]
Equivalently, for $n\ge 3$,
\[
\mathbf{H}_{n,\varepsilon}(x,y)=
\begin{pmatrix}
s+a^2 & -c & p \\
-\bar c & d & -c & \ddots \\
p & -\bar c & d & \ddots & \ddots \\
& \ddots & \ddots & \ddots & \ddots & \ddots \\
&& \ddots & \ddots & d & -c & p \\
&&& \ddots & -\bar c & d & -c \\
&&&& p & -\bar c & s+b^2
\end{pmatrix},
\]
with the obvious truncation when $n=2$, where
\[
\mathbf{H}_{2,\varepsilon}(x,y)=
\begin{pmatrix}
s+a^2 & -c \\
-\bar c & s+b^2
\end{pmatrix}.
\]
\end{proposition}

\begin{proof}
Expanding the product gives
\[
\mathbf{H}_{n,\varepsilon}(x,y)
=(\bar z\mathbf{I}-\mathbf{A}_n^{\mathsf{T}})(z\mathbf{I}-\mathbf{A}_n)-\varepsilon^2\mathbf{I}
=(|z|^2-\varepsilon^2)\mathbf{I}-\bar z\,\mathbf{A}_n-z\,\mathbf{A}_n^{\mathsf{T}}+\mathbf{A}_n^{\mathsf{T}}\mathbf{A}_n.
\]
Now $-\bar z\,\mathbf{A}_n-z\,\mathbf{A}_n^{\mathsf{T}}$ contributes the first off-diagonals
$-\bigl(az+b\bar z\bigr)=-c$ and $-\bigl(a\bar z+bz\bigr)=-\bar c$.
Moreover, $\mathbf{A}_n^{\mathsf{T}}\mathbf{A}_n$ has diagonal entries
$a^2,a^2+b^2,\dots,a^2+b^2,b^2$, has second off-diagonals equal to $ab=p$,
and has all first off-diagonal entries equal to $0$. The stated formula follows.
\end{proof}

\begin{corollary}\label{cor:Toeplitz-reduction}
Let $\Delta_0\coloneqq 1$, and for $m\ge 1$ let
\[
\Delta_m\coloneqq \det \mathbf{T}_m(d,c,p),
\]
where $\mathbf{T}_m(d,c,p)$ is the Hermitian Toeplitz pentadiagonal matrix
\[
\mathbf{T}_m(d,c,p)
=\begin{pmatrix}
d & -c & p \\
-\bar c & d & -c & \ddots \\
p & -\bar c & d & \ddots & \ddots \\
& \ddots & \ddots & \ddots & \ddots & \ddots \\
&& \ddots & \ddots & d & -c & p \\
&&& \ddots & -\bar c & d & -c \\
&&&& p & -\bar c & d
\end{pmatrix}.
\]
Then, for every $n\ge 2$,
\[
F_{n,\varepsilon}(x,y)
=\Delta_n-(a^2+b^2)\Delta_{n-1}+a^2b^2\Delta_{n-2}.
\]
\end{corollary}

\begin{proof}
By Proposition~\ref{prop:explicit-H},
\[
\mathbf{H}_{n,\varepsilon}(x,y)=\mathbf{T}_n(d,c,p)-b^2\mathbf{E}_{11}-a^2\mathbf{E}_{nn},
\]
where $\mathbf{E}_{jj}$ denotes the matrix with a single $1$ in the $(j,j)$ entry.
Since the determinant is affine in each of the two modified diagonal entries, we obtain
\[
\det \mathbf{H}_{n,\varepsilon}
=\det \mathbf{T}_n
-b^2\det \mathbf{T}_n^{(1)}
-a^2\det \mathbf{T}_n^{(n)}
+a^2b^2\det \mathbf{T}_n^{(1,n)},
\]
where $\mathbf{T}_n^{(1)}$ and $\mathbf{T}_n^{(n)}$ are obtained from $\mathbf{T}_n$ by deleting,
respectively, the first row and column and the last row and column, and
$\mathbf{T}_n^{(1,n)}$ is obtained by deleting both the first and the last row and column.
Because $\mathbf{T}_n$ is Toeplitz, these minors are precisely
$\mathbf{T}_{n-1}(d,c,p)$, $\mathbf{T}_{n-1}(d,c,p)$, and $\mathbf{T}_{n-2}(d,c,p)$.
Therefore
\[
F_{n,\varepsilon}(x,y)
=\Delta_n-(a^2+b^2)\Delta_{n-1}+a^2b^2\Delta_{n-2}.
\]
\end{proof}

\begin{remark}\label{rem:quartic-closed-form}
The Toeplitz determinants $\Delta_m$ admit a closed form in $m$. Let
$\rho_1,\rho_2,\rho_3,\rho_4$ be the roots of the quartic equation
\[
p\rho^4-c\rho^3+d\rho^2-\bar c\rho+p=0,
\]
that is,
\[
ab\rho^4-\bigl((a+b)x-i(b-a)y\bigr)\rho^3
+\bigl(x^2+y^2-\varepsilon^2+a^2+b^2\bigr)\rho^2
-\bigl((a+b)x+i(b-a)y\bigr)\rho
+ab=0.
\]
When the roots are distinct, one convenient closed form is
\[
\Delta_m
=p^m
\frac{
\det
\begin{pmatrix}
1&1&1&1\\
\rho_1&\rho_2&\rho_3&\rho_4\\
\rho_1^{m+2}&\rho_2^{m+2}&\rho_3^{m+2}&\rho_4^{m+2}\\
\rho_1^{m+3}&\rho_2^{m+3}&\rho_3^{m+3}&\rho_4^{m+3}
\end{pmatrix}
}{
\det
\begin{pmatrix}
1&1&1&1\\
\rho_1&\rho_2&\rho_3&\rho_4\\
\rho_1^2&\rho_2^2&\rho_3^2&\rho_4^2\\
\rho_1^3&\rho_2^3&\rho_3^3&\rho_4^3
\end{pmatrix}
},
\]
with the repeated-root case obtained by the usual confluent limit. Hence
\[
F_{n,\varepsilon}(x,y)
=\Delta_n-(a^2+b^2)\Delta_{n-1}+a^2b^2\Delta_{n-2}
\]
provides a closed form in $n$.
\end{remark}

\end{document}
